%%======%%======%0%======%%=======%%
\newpage
\section{Apêndices}
\label{sec:apendice}
%\addcontentsline{toc}{chapter}{Conclusão}

%%======%%======%0%======%%=======%%
\subsection{Uso de IA}
\begin{longtable}{>{\RaggedRight\arraybackslash}p{0.25\textwidth} >{\RaggedRight\arraybackslash}p{0.24\textwidth} >{\RaggedRight\arraybackslash}p{0.43\textwidth}} \caption{Participação da inteligência artificial generativa no desenvolvimento dos arquivos do site.} \label{tab:uso-ia-arquivos}\\ \hline \textbf{Arquivo do código} & \textbf{Função da IA} & \textbf{Análise geral crítica} \\ \hline \endfirsthead \hline \textbf{Arquivo do código} & \textbf{Função da IA} & \textbf{Análise geral crítica} \\ \hline \endhead \texttt{index.html} & \cellcolor{iaRefinamento}\textbf{Refinamento} & Revisão estrutural e semântica. A IA contribuiu para a organização da página inicial, acessibilidade e estrutura do \textit{landing page}, mas a composição final permaneceu sujeita a ajustes manuais. \\ \hline \texttt{login.html} & \cellcolor{iaRefinamento}\textbf{Refinamento} & Código explicitamente revisado e refinado com IA. A contribuição concentrou-se na estrutura da página, organização dos elementos e integração com os demais recursos do site. \\ \hline \texttt{maquinas.html} & \cellcolor{iaAcessibilidade}\textbf{Acessibilidade} & Uso de elementos semânticos, atributos ARIA e estrutura para diálogos. A IA auxiliou na implementação, mas a adequação ao funcionamento real exigiu integração com JavaScript e CSS. \\ \hline \texttt{membros.html} & \cellcolor{iaAcessibilidade}\textbf{Acessibilidade} & A estrutura emprega recursos de acessibilidade e organização semântica. Não há atribuição explícita de geração no arquivo, portanto a participação específica da IA não é quantificada. \\ \hline \texttt{repositorio.html} & \cellcolor{iaGeracao}\textbf{Geração} & Os filtros de pesquisa possuem auxílio explicitamente atribuído ao ChatGPT. A estrutura foi posteriormente integrada ao restante do sistema e às funções de persistência local. \\ \hline \texttt{adicionar-trabalho.html} & \cellcolor{iaGeracao}\textbf{Geração} & Trecho explicitamente declarado como gerado integralmente por IA. A tradução para o inglês e a adaptação ao funcionamento do projeto foram realizadas manualmente. \\ \hline \texttt{404.html} & \cellcolor{iaSemEvidencia}\textbf{Sem atribuição a IA} & Arquivo simples e predominantemente estrutural. O código disponível foi gerado integralmente por mãos humanas. \\ \hline \texttt{variables.css} & \cellcolor{iaRefinamento}\textbf{Refinamento} & O arquivo registra explicitamente o uso do Gemini Pro para definição e refinamento da abordagem visual, especialmente na organização das variáveis e implementação dos temas. \\ \hline \texttt{reset.css} & \cellcolor{iaSemEvidencia}\textbf{Sem atribuição} & Folha de estilo curta e estrutural. Sua função é normalizar o comportamento dos elementos antes da aplicação dos demais estilos. \\ \hline \texttt{base.css} & \cellcolor{iaCorrecao}\textbf{Correção} & O modo escuro foi explicitamente ajustado com auxílio do ChatGPT após problemas de implementação. A IA atuou principalmente na correção do comportamento visual. \\ \hline \texttt{layout.css} & \cellcolor{iaRefinamento}\textbf{Refinamento} & Responsável pela estrutura global, navegação, cabeçalho, rodapé e comportamento responsivo. O arquivo mostra forte integração com decisões de layout realizadas durante o desenvolvimento. \\ \hline \texttt{components.css} & \cellcolor{iaReestruturacao}\textbf{Reestruturação} & É o arquivo com maior participação de IA. GPT e Claude foram utilizados na elaboração e aprimoramento de grandes trechos, enquanto partes das máquinas e membros foram desenvolvidas manualmente. O próprio código registra dependência significativa da IA. \\ \hline \texttt{adicionar-trabalho.js} & \cellcolor{iaCorrecao}\textbf{Correção} & O código partiu de uma tentativa própria que não funcionava adequadamente e foi posteriormente adaptado com auxílio do GPT. A IA atuou principalmente na correção e adequação da lógica de cadastro. \\ \hline \texttt{auth.js} & \cellcolor{iaReestruturacao}\textbf{Reestruturação} & A lógica de autenticação e persistência da sessão foi reestruturada com auxílio do GPT. O uso de \texttt{localStorage} resolveu problemas de perda de sessão durante navegação e atualização das páginas. \\ \hline \texttt{carrossel.js} & \cellcolor{iaAcessibilidade}\textbf{Acessibilidade} & O GPT auxiliou na lógica de carregamento e navegação do carrossel, enquanto o Gemini Pro foi utilizado no tratamento relacionado à acessibilidade e preferência por redução de movimento. \\ \hline \texttt{data.js} & \cellcolor{iaSemEvidencia}\textbf{Sem atribuição significativa} & Arquivo predominantemente voltado ao armazenamento dos dados iniciais do protótipo. A participação específica da IA na sua elaboração não é relevante, usada apenas para validação. \\ \hline \texttt{i18n.js} & \cellcolor{iaAcessibilidade}\textbf{Internacionalização} & As estruturas de tradução foram geradas com auxílio do GPT, enquanto a lógica de seleção de idioma foi posteriormente reestruturada. Representa uma das maiores contribuições da IA para a expansão funcional do projeto. \\ \hline \texttt{login-carrossel.js} & \cellcolor{iaRefinamento}\textbf{Refinamento} & Responsável pelo comportamento visual do carrossel da página de login. O arquivo apresenta lógica simples e integrada ao CSS, com tratamento da preferência por redução de movimento. \\ \hline \texttt{maquinas.js} & \cellcolor{iaReestruturacao}\textbf{Reestruturação} & O código registra explicitamente o uso de IA para reestruturar a apresentação das imagens dos equipamentos e reduzir redundâncias. A contribuição ocorreu sobre uma lógica já em desenvolvimento. \\ \hline \texttt{membros.js} & \cellcolor{iaRefinamento}\textbf{Refinamento} & Implementa renderização dinâmica, categorias, links acadêmicos e abas. A estrutura é fortemente integrada ao sistema de tradução e aos dados locais, embora não haja atribuição significativa de autoria da IA. \\ \hline \texttt{repositorio.js} & \cellcolor{iaCorrecao}\textbf{Correção} & A IA foi utilizada para fazer o mecanismo de exibição e filtragem funcionar. O código também incorpora \texttt{escapeHtml}, demonstrando preocupação com tratamento seguro do conteúdo inserido no DOM. \\ \hline \texttt{tema.js} & \cellcolor{iaRefinamento}\textbf{Refinamento} & Implementa a alternância entre temas e sua persistência no navegador. A lógica é compacta e integrada às variáveis CSS, permitindo separar a aparência da lógica de controle. \\ \hline \texttt{ui.js} & \cellcolor{iaRefinamento}\textbf{Refinamento} & Controla menu responsivo, estado de expansão e atualização do ano do rodapé. O código demonstra integração entre comportamento visual, navegação e tradução. \\ \hline \texttt{video.js} & \cellcolor{iaRefinamento}\textbf{Refinamento} & Controla reprodução, volume, pausa, reinício e reprodução condicionada à visibilidade do vídeo. A lógica utiliza recursos nativos do navegador e permanece relativamente independente. \\ \hline \end{longtable} \vspace{3mm} \noindent \textbf{Legenda da classificação:} \quad \cellcolor{iaGeracao}\textbf{Geração} \quad \cellcolor{iaCorrecao}\textbf{Correção} \quad \cellcolor{iaRefinamento}\textbf{Refinamento} \quad \cellcolor{iaAcessibilidade}\textbf{Acessibilidade/Internacionalização} \quad \cellcolor{iaReestruturacao}\textbf{Reestruturação} \quad \cellcolor{iaSemEvidencia}\textbf{Sem atribuição}

\subsection{metricas}
\subsubsection{W3C CSS}

\begin{figure}[H]
	\centering
	
	\begin{subfigure}[b]{0.48\textwidth}
		\centering
		\includegraphics[width=\linewidth]{"../2-imagens/W3C CSS/base"}
		\caption{Base}
		\label{fig:w3c-base}
	\end{subfigure}
	\hfill
	\begin{subfigure}[b]{0.48\textwidth}
		\centering
		\includegraphics[width=\linewidth]{"../2-imagens/W3C CSS/components"}
		\caption{Components}
		\label{fig:w3c-components}
	\end{subfigure}
	
	\vspace{5mm}
	
	\begin{subfigure}[b]{0.48\textwidth}
		\centering
		\includegraphics[width=\linewidth]{"../2-imagens/W3C CSS/layout"}
		\caption{Layout}
		\label{fig:w3c-layout}
	\end{subfigure}
	\hfill
	\begin{subfigure}[b]{0.48\textwidth}
		\centering
		\includegraphics[width=\linewidth]{"../2-imagens/W3C CSS/reset"}
		\caption{Reset}
		\label{fig:w3c-reset}
	\end{subfigure}
	
	\vspace{5mm}
	
	\begin{subfigure}[b]{0.48\textwidth}
		\centering
		\includegraphics[width=\linewidth]{"../2-imagens/W3C CSS/variables"}
		\caption{Variables}
		\label{fig:w3c-variables}
	\end{subfigure}
	
	\caption{Resultados da validação dos arquivos CSS utilizando o validador W3C.}
	\label{fig:w3c-css}
\end{figure}

\subsubsection{W3C HTML}

\begin{figure}[H]
	\centering
	
	\begin{subfigure}[b]{0.48\textwidth}
		\centering
		\includegraphics[width=\linewidth]{"2-imagens/W3C HTML/login/login"}
		\caption{Login}
		\label{fig:w3c-html-login}
	\end{subfigure}
	\hfill
	\begin{subfigure}[b]{0.48\textwidth}
		\centering
		\includegraphics[width=\linewidth]{"2-imagens/W3C HTML/maq/maquinas errors"}
		\caption{Máquinas}
		\label{fig:w3c-html-maquinas}
	\end{subfigure}
	
	\vspace{5mm}
	
	\begin{subfigure}[b]{0.48\textwidth}
		\centering
		\includegraphics[width=\linewidth]{"2-imagens/W3C HTML/memb/membros .png"}
		\caption{Membros}
		\label{fig:w3c-html-membros}
	\end{subfigure}
	\hfill
	\begin{subfigure}[b]{0.48\textwidth}
		\centering
		\includegraphics[width=\linewidth]{"2-imagens/W3C HTML/reposit/repositorio"}
		\caption{Repositório}
		\label{fig:w3c-html-repositorio}
	\end{subfigure}
	
	\caption{Resultados da validação dos arquivos HTML utilizando o validador W3C.}
	\label{fig:w3c-html}
\end{figure}

\subsubsection{webAIM}

\begin{figure}[H]
	\centering
	
	\begin{subfigure}[b]{0.48\textwidth}
		\centering
		\includegraphics[width=\linewidth]{"2-imagens/webAIM/TEMA CLARO/cards"}
		\caption{Cards}
		\label{fig:webaim-claro-cards}
	\end{subfigure}
	\hfill
	\begin{subfigure}[b]{0.48\textwidth}
		\centering
		\includegraphics[width=\linewidth]{"2-imagens/webAIM/TEMA CLARO/texto principal"}
		\caption{Texto principal}
		\label{fig:webaim-claro-principal}
	\end{subfigure}
	
	\vspace{5mm}
	
	\begin{subfigure}[b]{0.48\textwidth}
		\centering
		\includegraphics[width=\linewidth]{"2-imagens/webAIM/TEMA CLARO/texto secundario"}
		\caption{Texto secundário}
		\label{fig:webaim-claro-secundario}
	\end{subfigure}
	\hfill
	\begin{subfigure}[b]{0.48\textwidth}
		\centering
		\includegraphics[width=\linewidth]{"2-imagens/webAIM/TEMA ESCURO/texto principal"}
		\caption{Texto principal}
		\label{fig:webaim-escuro-principal}
	\end{subfigure}
	
	\vspace{5mm}
	
	\begin{subfigure}[b]{0.48\textwidth}
		\centering
		\includegraphics[width=\linewidth]{"2-imagens/webAIM/TEMA ESCURO/texto secundario"}
		\caption{Texto secundário}
		\label{fig:webaim-escuro-secundario}
	\end{subfigure}
	\hfill
	\begin{subfigure}[b]{0.48\textwidth}
		\centering
		\includegraphics[width=\linewidth]{"2-imagens/webAIM/header selecionado"}
		\caption{Cabeçalho}
		\label{fig:webaim-header}
	\end{subfigure}
	
	\caption{Testes de contraste visual realizados com a ferramenta WebAIM para os diferentes elementos e temas da interface.}
	\label{fig:webaim}
\end{figure}

\subsubsection{LightHouse}

\begin{figure}[H]
	\centering
	
	\begin{subfigure}[b]{0.31\textwidth}
		\centering
		\includegraphics[width=\linewidth]{"2-imagens/home lighthouse"}
		\caption{Página inicial.}
		\label{fig:home-lighthouse}
	\end{subfigure}
	\hfill
	\begin{subfigure}[b]{0.31\textwidth}
		\centering
		\includegraphics[width=\linewidth]{"2-imagens/login lighthouse"}
		\caption{Login.}
		\label{fig:login-lighthouse}
	\end{subfigure}
	\hfill
	\begin{subfigure}[b]{0.31\textwidth}
		\centering
		\includegraphics[width=\linewidth]{"2-imagens/maq lighthouse"}
		\caption{Máquinas.}
		\label{fig:maq-lighthouse}
	\end{subfigure}
	
	\vspace{5mm}
	
	\begin{subfigure}[b]{0.31\textwidth}
		\centering
		\includegraphics[width=\linewidth]{"2-imagens/membros lighthouse"}
		\caption{Membros.}
		\label{fig:membros-lighthouse}
	\end{subfigure}
	\hfill
	\begin{subfigure}[b]{0.31\textwidth}
		\centering
		\includegraphics[width=\linewidth]{"2-imagens/repositorio lighthouse"}
		\caption{Repositório.}
		\label{fig:repositorio-lighthouse}
	\end{subfigure}
	
	\caption{Resultados das análises realizadas pelo Lighthouse nas páginas do sistema.}
	\label{fig:lighthouse-resultados}
\end{figure}

\subsection{Check list Boas práticas}

\begin{longtable}{p{3.2cm} p{7.2cm} p{3.0cm}}
	
	\caption{Boas práticas utilizadas no desenvolvimento do site.}
	\label{tab:boas-praticas}
	\\
	
	\toprule
	\textbf{Prática} &
	\textbf{Implementação} &
	\textbf{Arquivo} \\
	\midrule
	\endfirsthead
	
	\toprule
	\textbf{Prática} &
	\textbf{Implementação} &
	\textbf{Arquivo} \\
	\midrule
	\endhead
	
	\midrule
	\multicolumn{3}{r}{Continua na próxima página} \\
	\endfoot
	
	\bottomrule
	\endlastfoot
	
	DOCTYPE HTML5 &
	Utilização da declaração \texttt{doctype} correspondente ao HTML5. &
	\texttt{index.html}, linha 2 \\
	
	Atributo \texttt{lang} &
	Definição do idioma principal da página como português do Brasil por meio do atributo \texttt{lang}. &
	\texttt{index.html}, linha 3 \\
	
	Meta viewport &
	Configuração da área de visualização para adaptação do conteúdo às diferentes dimensões de tela. &
	\texttt{index.html}, linha 6 \\
	
	CSS externo &
	Separação dos estilos em arquivos específicos para reset, variáveis, estilos básicos, layout e componentes. &
	\texttt{index.html}, linhas 21--25 \\
	
	Unidades relativas &
	Utilização de unidades e funções relativas, como \texttt{rem}, porcentagem, \texttt{vw}, \texttt{ch}, \texttt{min()} e \texttt{clamp()}. &
	\texttt{css/}, diversos arquivos \\
	
	Imagens com \texttt{alt} &
	As imagens utilizadas no site possuem o atributo alternativo \texttt{alt}. &
	\texttt{index.html} \\
	
	Dimensões explícitas de imagens &
	As imagens do carrossel possuem dimensões definidas por meio dos atributos \texttt{width} e \texttt{height}. &
	\texttt{index.html}, linhas 116--140 \\
	
	Flexbox e Grid &
	Utilização de Flexbox e CSS Grid para organização dos elementos e estrutura das páginas. &
	\texttt{css/layout.css} \\
	
	Ausência de estilos inline &
	Os estilos são mantidos nos arquivos CSS externos, sem utilização de atributos \texttt{style} diretamente nos elementos HTML. &
	Arquivos HTML e CSS \\
	
	IDs não duplicados &
	Os identificadores utilizados nos elementos das páginas são definidos individualmente. &
	Arquivos HTML \\
	
	Links internos &
	Utilização de links para navegação entre as diferentes páginas que compõem o site. &
	Arquivos HTML \\
	
	Foco visível &
	Utilização de \texttt{:focus-visible} para apresentar uma indicação visual durante a navegação por teclado. &
	\texttt{css/base.css}, linhas 48--52 \\
	
	Contraste &
	Utilização de cores centralizadas nas variáveis CSS para padronização da interface nos diferentes temas. &
	\texttt{css/variables.css} \\
	
	Labels associados aos campos &
	Os campos dos formulários possuem elementos \texttt{label} associados aos respectivos campos por meio dos atributos \texttt{for} e \texttt{id}. &
	\texttt{login.html} e formulários \\
	
	Navegação por teclado &
	Implementação de elementos navegáveis por teclado, incluindo botão de menu, link de salto e indicação visual de foco. &
	\texttt{index.html} e \texttt{css/base.css} \\
	
	Estrutura semântica &
	Utilização de elementos semânticos do HTML, como \texttt{header}, \texttt{nav}, \texttt{main}, \texttt{section} e \texttt{footer}. &
	Arquivos HTML \\
	
	Variáveis CSS &
	Definição de valores reutilizáveis dentro de \texttt{:root}, permitindo centralizar cores e demais propriedades da interface. &
	\texttt{css/variables.css} \\
	
\end{longtable}

\subsection{Histórico git}

\begin{lstlisting}[style=gitstyle,
	caption={Histórico de commits do desenvolvimento do projeto.},
	label={lst:historico-git}]
	* fa5c243 (HEAD -> main, origin/main, origin/HEAD) docs: escrita do relatorio
	* 8872e58 perf: compressao de outras fotos
	* 07e8647 perf: compressao de videos e imagens para aumentar lighthouse
	* fa3e505 Remove métricas do repositório
	* bb9de45 docs: atualiza e adiciona um README V2
	* 1f99cb4 fix: correcao de imagem do quatro pontas
	* dcc5abc style: adicao de um favIcon
	* 9f48717 style: ajustes e melhorias nos botoes clicavies de contato com os membros
	* 64f1b3d feat: adicao de mailto nos membros e alteracao do footer
	* 9cbc549 feat: adiciona comentario sobre as maquinas
	* 04ae87d feat: adiciona campo de nome, inst. e motivacao para agendamentos
	* fd01438 perf: Melhora de perfomance adicionando o load lazy
	* 6d99d60 perf: compressao de fotos e videos - melhorar a performance do lighthouse
	* dd0986d Configura GitHub Pages
	* 5b75d3f ajeitando pjt nas pasta Site E Relatorio
	* 3cee70b mexendo no gitignore p arquivos latex
	* 726c986 update
	* 1e32958 style: Adicao de transicao entre idiomas e temas (claro e escuro) e adicao do logo LIFS
	* 5737548 feat: adicao de video institucional e mapa de localizaçao do lab
	* 72a30da fix:Correçao do auth.js que perdia o status de login + Ajuste de membros nao atualizar o idioma
	* 2dd885d feat: Traduçao das paginas de membros e maquinas V2
	* 460383b feat: Traduçao da pagina de erro 404
	* a23b6a9 feat: Traduçao da pagina de adicionar trabalho
	* e4253ed feat: Traduçao da pagina de login
	* 70481d7 feat: Traduçao da pagina de membros
	* e1e16aa feat: Traduçao das paginas de maquinas e repositorio
	* 678f665 feat: adiçao da traduçao em ingles da pagina HOME/INICIO, para acessibilidade de internacionalizandos.
	* f7c428a chore: cria a pasta js/i18n.js para tradução do site em ingles
	* 353e71b feat: adicao do tema ESCURO e CLARO no site
	* 894860b style: Adicao do carrossel na landingpage no inicio.
	* c2d018a style:adicao de um carrossel de fundo na tela de login. alteracoes em java css e html basicas.
	* 368b76b style: Mudanças relevantes no estilo de membros e maquinas, bem como nas fontes principais
	* b934165 chore: cria a pasta para add trabalho .js
	* 5b86a34 style: correcao de ortografia nas legendas
	* dae911f style:ajuste para imagens serem clicaveis
	* 8d2253d style:imagem das maquinas estilizada
	* 9f99d0a style:Adicao das imagens de membros e maquinas
	* c36fff7 refactor: Codigos HTML refinados por IA e maquinas.js adicionadas
	* 7b10636 feat: ajuste das partes faltantes
	* b3cae9b feat: adiciona pagina e logica de autenticacao de usuarios
	* 93a8197 Feature: Adicao do carrosel de imagens
	* 4561067 STYLE: Adicao de trechos do CSS para layout
	* 5ed7317 STYLE: Adicao de trechos do CSS
	* 097b0d6 feat: estrutura inicial do index.html
	* 8e1a315 docs: add de um pseudo readme
	* daad2b1 chore: cria estrutura de pastas, arquivos vazios e gitignore
	* d111ce8 chore: estrutura inicial do projeto
\end{lstlisting}